\section{Arithmetic models and projections}\label{sec:notation}

All polynomials are formal and commutative. A circuit over \(F\) uses
binary addition, subtraction, and multiplication gates. Its size \(s\)
counts the operation gates. A formula is a circuit with a tree of
gate occurrences. We allow an arbitrary affine form in the input
variables at each leaf, free of charge. An ABP sums products of affine edge labels along source--sink paths in
a directed acyclic graph. Its size counts both vertices and edges. These input conventions preserve the
polynomial-size classes for families of polynomially bounded arity.
We impose no bound on the degree of intermediate computations.
Throughout, logarithms have base two.

Series and parallel composition turn a formula into a branching program
of linear size. A circuit can evaluate a branching program in
topological order, summing at each vertex the value at each predecessor
times the connecting edge label. Finally, a family in \(\VNP\) has the
form
\[
          f_n(X)=\sum_{b\in\{0,1\}^{p(n)}}g_n(b,X)
\]
for a polynomially bounded \(p(n)\) and \((g_n)\in\VP\). Taking
\(p(n)=0\) gives the last inclusion in Theorem~\ref{thm:main}.

\paragraph{Degree-component border size.}
Fix a homogeneous target of degree \(D\) and put \(K=\overline F\).
The target has \emph{degree-component border size} at most \(s\) if
its coefficient vector lies in the Zariski closure, over \(K\), of the
degree-\(D\) components of outputs of size-\(s\) or smaller computations
in the specified model. Approximating outputs can have terms in other
degrees. The following lemma accounts for those terms when a substitution
uses constants and lowers the degree.

\begin{lemma}[Projection]\label{lem:projection}
Let \(f\) be homogeneous of degree \(D\) and suppose that substituting
affine forms for its variables gives a homogeneous polynomial \(g\).
If \(f\) has degree-component border size at most \(s\), then in the
same model \(g\) has degree-component border size at most
\(O((s+1)(D+1)^2)\) for circuits and \(O((s+1)^2)\) for formulas or ABPs.
\end{lemma}
\begin{proof}
First extract the degree-\(D\) component of each approximating output.
A circuit computes all components through degree \(D\) at each gate,
using convolution at multiplication gates. This uses
\(O((s+1)(D+1)^2)\) gates. A size-\(s\) formula or ABP has output degree
at most \(s+1\). Over the infinite field \(K\), interpolation in
\(t\) at \(s+2\) distinct values isolates degree \(D\) in the scaled
output \(f_{\rm out}(tX)\). If \(D>s+1\) the component is zero.
Combining the scaled copies otherwise costs \(O((s+1)^2)\).

Substitute the prescribed affine forms in the extracted component
and read its component of degree \(\deg g\). For the target this
returns \(g\). From the degree-\(D\) coefficient space to the
degree-\(\deg g\) coefficient space, the resulting map is linear.
It therefore takes the original target into the closure of the
transformed outputs, with the stated size bound.
\end{proof}

\paragraph{Subgraphs and smaller alphabets.}
The polynomial \(Z_{G,q}\) is homogeneous of degree \(|V(G)|\) and
set-multilinear in the vertex tables. For a subgraph \(H\subseteq G\)
and \(2\le q'\le q\), a substitution of variables, zero, and one sends
\(Z_{G,q}\) to \(Z_{H,q'}\). Pin each unused edge to state one and allow
only the first \(q'\) states on the retained edges. A table entry
inconsistent with these choices is set to zero. A consistent entry
becomes the corresponding variable at a retained vertex and becomes
one at a deleted vertex. Each labeling of \(H\) has exactly one
surviving extension to \(G\). This gives the projections used in the
three witness proofs. Lemma~\ref{lem:graph-minor} extends the construction
to graph minors.
